% !TeX program = xelatex
% Compile với XeLaTeX: xelatex proposal.tex && xelatex proposal.tex
\documentclass[a4paper,11pt]{article}

\usepackage{fontspec}
\setmainfont{Noto Serif}
\setsansfont{Noto Sans}
\usepackage{polyglossia}
\setmainlanguage{vietnamese}

\usepackage[a4paper,margin=2.5cm]{geometry}
\usepackage{amsmath,amssymb}
\usepackage{graphicx}
\usepackage{booktabs}
\usepackage{longtable}
\usepackage{array}
\usepackage{enumitem}
\usepackage{xcolor}
\usepackage{hyperref}
\hypersetup{
    colorlinks=true,
    linkcolor=blue!50!black,
    urlcolor=blue!50!black,
    citecolor=blue!50!black,
    pdftitle={Đề xuất dự án: Digital Twin cho SADGS},
    pdfauthor={Digital Twin GS}
}
\usepackage{fancyhdr}
\usepackage{titlesec}
\usepackage{tcolorbox}
\tcbuselibrary{skins,breakable}

\titleformat{\section}{\normalfont\Large\bfseries}{\thesection.}{0.5em}{}
\titleformat{\subsection}{\normalfont\large\bfseries}{\thesubsection.}{0.5em}{}

\pagestyle{fancy}
\fancyhf{}
\fancyhead[L]{\small Đề xuất dự án -- Digital Twin cho SADGS}
\fancyhead[R]{\small\thepage}
\renewcommand{\headrulewidth}{0.4pt}

\newtcolorbox{luuy}[1][]{colback=yellow!8,colframe=orange!60!black,fonttitle=\bfseries,title=Lưu ý trung thực,boxrule=0.6pt,#1}
\newtcolorbox{muctieu}[1][]{colback=blue!5,colframe=blue!40!black,fonttitle=\bfseries,#1}

\title{\textbf{Đề xuất Dự án}\\[4pt]
\Large Digital Twin cho 3D Gaussian Splatting:\\
Nghiên cứu, Xác minh và Tài liệu hoá\\
Structure-Aware Densification (SADGS)}
\author{Digital Twin GS}
\date{\today}

\begin{document}

\maketitle
\thispagestyle{empty}

\begin{abstract}
\noindent
Tài liệu này đề xuất một dự án học thuật nhằm xây dựng một bộ tài liệu ``digital twin'' đầy
đủ -- sách kỹ thuật, slide thuyết trình, và hình minh hoạ được sinh từ số liệu/thực nghiệm
thật -- cho phương pháp \emph{Structure-Aware Densification for Gaussian Splatting} (SADGS,
Lyu et al., SIGGRAPH 2026). Điểm khác biệt cốt lõi của dự án so với việc chép lại paper là:
\textbf{mọi công thức, tên hàm, số dòng code và tham số mặc định trình bày trong tài liệu đều
phải được đối chiếu (grep/read) trực tiếp với mã nguồn thật trong \texttt{SADGS/}}, và mọi
hình matplotlib minh hoạ phải được tính từ dữ liệu/ảnh/checkpoint thật thay vì số liệu bịa.
Báo cáo này trình bày bối cảnh, mục tiêu, phạm vi, phương pháp luận, kế hoạch triển khai và
các hạn chế đã biết của dự án tại thời điểm đề xuất.
\end{abstract}

\vspace{0.5cm}
\tableofcontents

\newpage

\section{Bối cảnh và động lực}

\subsection{Bài toán gốc}
3D Gaussian Splatting (3DGS, Kerbl et al. 2023) biểu diễn cảnh 3D bằng một tập hợp Gaussian
dị hướng, được khởi tạo từ point cloud SfM (COLMAP) và tối ưu bằng render khả vi. Một khâu
quan trọng của pipeline là \emph{Adaptive Density Control} (ADC): quyết định khi nào một
Gaussian cần được nhân đôi (clone), tách (split), hay loại bỏ (prune) dựa trên gradient vị trí
tích luỹ qua các bước huấn luyện.

\subsection{Hạn chế của ADC nguyên bản và đóng góp của SADGS}
ADC nguyên bản chỉ dùng độ lớn gradient vị trí trung bình -- một tín hiệu gián tiếp, dễ bị
nhiễu bởi hiện tượng \emph{gradient cancellation} khi nhiều view ``triệt tiêu'' lẫn nhau.
SADGS thay thế bằng một tín hiệu trực tiếp hơn: so khớp \emph{screen-space extent} của mỗi
Gaussian với \emph{cấu trúc texture đa tỉ lệ} (multiscale structure tensor) của ảnh mục tiêu,
từ đó suy ra tỉ số $\eta$ và thực hiện \emph{split dị hướng} (anisotropic split) thay vì
clone/split đẳng hướng, cộng thêm bước \emph{multiview-consistency} để tránh quyết định dựa
trên một view nhiễu đơn lẻ. Các hàm lõi nằm tại \texttt{SADGS/scene/gaussian\_model.py}:
\texttt{densify\_and\_split\_structgs}, \texttt{densify\_and\_clone\_structgs},
\texttt{densify\_and\_prune\_structgs}, \texttt{final\_prune\_structgs}.

\subsection{Vấn đề đã phát hiện trong tài liệu hiện có}
Trước khi đề xuất này được viết, một đợt rà soát toàn diện (đối chiếu từng chương sách, từng
slide, từng script matplotlib với mã nguồn thật trong \texttt{SADGS/}) đã phát hiện một lượng
đáng kể nội dung \textbf{không khớp với code đang chạy thật}, cụ thể:
\begin{itemize}[leftmargin=1.2em]
    \item Một cơ chế ``FastGS'' được mô tả như một baseline đã được triển khai và kiểm thử
    (hàm \texttt{compute\_gaussian\_score\_fastgs}, \texttt{render\_fastgs}, framework 7 bước)
    -- trong khi thực tế \texttt{SADGS/} chỉ có \emph{một dòng comment} và một dòng credit
    trong README; các hàm này không tồn tại.
    \item Một điểm số ``Importance'' dựa trên lỗi pixel trong footprint (\texttt{metric\_map})
    được trình bày như cơ chế quyết định prune/clone thật -- trong khi lệnh gọi thật trong
    \texttt{train.py} truyền \texttt{importance\_score=accum\_view\_count} (một bộ đếm số
    view nhìn thấy, hoàn toàn khác về ngữ nghĩa).
    \item \texttt{final\_prune\_structgs} và \texttt{expand\_undersized\_gs} bị comment out
    trong \texttt{train.py} -- tức là \textbf{không chạy trong pipeline mặc định} -- nhưng
    nhiều đoạn tài liệu mô tả như thể chúng đang hoạt động.
    \item Một số chương trích dẫn một package \texttt{pipeline/} (\texttt{trainer.py},
    \texttt{config.py}, v.v.) như thể nó là kiến trúc thật của repo này, trong khi đó thực ra
    là cấu trúc của một fork Colab riêng (``SADGS-lite''), không tồn tại trong
    \texttt{SADGS/}.
    \item Nhiều hình matplotlib minh hoạ dùng số liệu \texttt{np.random}/tự bịa, được gắn nhãn
    ngầm như thể là kết quả đo thật.
\end{itemize}

\begin{luuy}
Đợt rà soát này đã sửa phần lớn các điểm trên (xoá nội dung bịa hoàn toàn, đính chính các nội
dung có công thức thật nhưng bị mô tả sai ngữ cảnh live/dead-code, tạo lại hình minh hoạ từ dữ
liệu/ảnh/checkpoint thật). Dự án được đề xuất trong tài liệu này chính thức hoá quy trình xác
minh đó thành phương pháp luận bắt buộc cho mọi nội dung được thêm mới về sau, thay vì xử lý
một lần rồi thôi.
\end{luuy}

\section{Mục tiêu dự án}

\begin{muctieu}
\textbf{Mục tiêu tổng quát:} Xây dựng một bộ tài liệu digital-twin cho SADGS mà mọi khẳng
định kỹ thuật đều \emph{truy vết được} (traceable) tới một trong hai nguồn: (a) mã nguồn thật
trong \texttt{SADGS/}, hoặc (b) dữ liệu/ảnh/log thật được lưu trong repo -- không có nội dung
nào được trình bày như sự thật mà không kiểm chứng được.
\end{muctieu}

\noindent Mục tiêu cụ thể:
\begin{enumerate}[leftmargin=1.4em]
    \item \textbf{Giải thích chính xác pipeline 3DGS gốc} (khởi tạo, rasterizer khả vi, loss,
    backprop, ADC) làm nền tảng trước khi trình bày phần mở rộng SADGS.
    \item \textbf{Giải thích chính xác cơ chế SADGS}: structure tensor đa tỉ lệ, tỉ số $\eta$,
    split dị hướng, multiview-consistency pruning -- kèm chỉ rõ những phần đã cài đặt nhưng
    \emph{không active} trong pipeline mặc định.
    \item \textbf{Sinh hình minh hoạ từ dữ liệu thật}: mọi biểu đồ matplotlib phải được tính
    từ ảnh huấn luyện thật (\texttt{DOCS/assets/}), checkpoint Gaussian thật (\texttt{.ply}),
    hoặc log CSV thật (\texttt{history.csv}, \texttt{leaderboard.csv}) -- không dùng số liệu
    tự bịa để minh hoạ một phép đo.
    \item \textbf{Chạy lại benchmark chính thức cho SADGS thật}. Số liệu hiện có trong
    \texttt{DOCS/README.md} là của pipeline \emph{trước khi} tích hợp structure-aware
    densification (``FastGS-lite''); cần một lần chạy \texttt{SADGS/run\_train.sh} trên cùng
    scene để có số liệu SADGS chính thức, so sánh công bằng với baseline.
    \item \textbf{Duy trì quy trình kiểm chứng liên tục}: mọi lần cập nhật tài liệu sau này
    đều đi kèm bước đối chiếu lại với code, tránh tài liệu trôi dạt (drift) khỏi thực tế như
    đã xảy ra trước đợt rà soát.
\end{enumerate}

\section{Phạm vi dự án}

\subsection{Trong phạm vi}
\begin{itemize}[leftmargin=1.2em]
    \item Sách kỹ thuật \texttt{DOCS/BOOK/} (18 chương + phụ lục lời giải), tiếng Việt.
    \item Hai bộ slide LaTeX/Beamer: bản đầy đủ (\texttt{Slides67/main.tex}, tra cứu sâu) và
    bản 45 phút (\texttt{Slides67/Slides45min/main45.tex}, báo cáo chính).
    \item Toàn bộ script sinh hình \texttt{matplotlib} kèm hình \texttt{.png} tương ứng.
    \item Một lượt chạy benchmark chính thức cho SADGS thật trên hạ tầng Colab T4.
    \item Quy trình kiểm chứng (checklist đối chiếu code) áp dụng cho mọi nội dung mới.
\end{itemize}

\subsection{Ngoài phạm vi}
\begin{itemize}[leftmargin=1.2em]
    \item Sửa đổi/mở rộng thuật toán SADGS gốc (dự án chỉ tài liệu hoá, không đề xuất thuật
    toán mới).
    \item Tối ưu hiệu năng CUDA của rasterizer (\texttt{SADGS/submodules/}).
    \item Chạy benchmark trên nhiều bộ dữ liệu chuẩn (Mip-NeRF360, Tanks\&Temples) ngoài scene
    của cuộc thi \texttt{VAI\_NVS\_DATA\_ROUND2} -- có thể là hướng mở rộng ở giai đoạn sau.
\end{itemize}

\section{Phương pháp luận}

\subsection{Nguyên tắc xác minh bắt buộc}
Mọi khẳng định kỹ thuật thêm vào tài liệu phải đi qua quy trình sau trước khi được chấp nhận:
\begin{enumerate}[leftmargin=1.4em]
    \item \textbf{Grep/Read mã nguồn thật}: tên hàm, đường dẫn file, số dòng, giá trị tham số
    mặc định phải được tìm thấy nguyên văn trong \texttt{SADGS/}.
    \item \textbf{Phân biệt live path vs.\ dead code}: nếu một hàm tồn tại trong code nhưng lời
    gọi của nó bị comment out (như \texttt{final\_prune\_structgs}), tài liệu phải nói rõ điều
    này, không được ngầm định là đang chạy.
    \item \textbf{Dữ liệu minh hoạ phải thật hoặc dán nhãn rõ}: nếu một hình cần số liệu để
    minh hoạ một công thức, số liệu đó phải đến từ (a) chạy công thức thật trên ảnh/checkpoint
    thật, hoặc (b) một ví dụ đồ chơi (toy) được dán nhãn tường minh là ``minh hoạ, không phải
    số đo thật'' ngay trong tiêu đề/chú thích hình.
    \item \textbf{Không suy diễn thay cho code}: khi không chắc chắn, tài liệu ghi rõ
    \texttt{[cần xác minh lại]} thay vì khẳng định.
\end{enumerate}

\subsection{Quy trình sinh hình minh hoạ}
\begin{enumerate}[leftmargin=1.4em]
    \item Đọc công thức thật trong \texttt{SADGS/} (ví dụ: \texttt{get\_structure\_tensor\_torch},
    \texttt{densify\_and\_split\_structgs}).
    \item Tái hiện công thức đó bằng \texttt{numpy} thuần (môi trường viết tài liệu không có
    \texttt{torch}), có trích dẫn \texttt{file:line} ngay trong docstring của script.
    \item Chạy trên dữ liệu thật: ảnh huấn luyện (\texttt{DOCS/assets/samples\_*.png}),
    checkpoint Gaussian thật (\texttt{point\_cloud\_iter5000.ply}), hoặc log CSV thật
    (\texttt{history.csv}, \texttt{leaderboard.csv}).
    \item Lưu hình với tên mô tả, chạy thử script để xác nhận không lỗi trước khi đưa vào sách.
\end{enumerate}

\subsection{Công cụ}
Biên soạn bằng Markdown (sách) và LaTeX/Beamer + XeLaTeX (slide, do cần font Unicode tiếng
Việt qua \texttt{fontspec}/\texttt{polyglossia}). Hình minh hoạ bằng \texttt{matplotlib}. Mã
nguồn tham chiếu: \texttt{SADGS/} (Python/PyTorch/CUDA).

\section{Kế hoạch triển khai và sản phẩm bàn giao}

\begin{longtable}{@{}>{\raggedright\arraybackslash}p{3.6cm} p{9.5cm}@{}}
\toprule
\textbf{Giai đoạn} & \textbf{Nội dung} \\
\midrule
\endhead
G0 -- Rà soát nền & Đối chiếu toàn bộ sách/slide/hình hiện có với mã nguồn thật; xoá nội dung
bịa; sửa trích dẫn sai; tạo lại hình từ dữ liệu thật. \emph{(Đã hoàn thành phần lớn tại thời
điểm đề xuất -- xem Mục~6.)} \\
G1 -- Chuẩn hoá quy trình & Viết checklist xác minh (Mục~4.1) thành quy tắc bắt buộc cho mọi
đóng góp nội dung mới; thiết lập chỗ lưu trữ số liệu thật (\texttt{DOCS/assets/}) làm nguồn
tham chiếu duy nhất cho hình minh hoạ. \\
G2 -- Chạy benchmark SADGS chính thức & Chạy \texttt{SADGS/run\_train.sh} đầy đủ 30\,000
iteration trên scene \texttt{HCM0539} (hoặc scene chuẩn khác), ghi \texttt{history.csv}/
\texttt{leaderboard.csv} thật cho SADGS (khác với số liệu FastGS-lite hiện có). \\
G3 -- Cập nhật so sánh định lượng & Thay số liệu ``tham khảo'' hiện tại trong
\texttt{DOCS/README.md} và chương 17 bằng số liệu SADGS thật; cập nhật slide so sánh
3DGS gốc / SADGS tương ứng. \\
G4 -- Hoàn thiện \& build cuối & Build PDF cả hai bộ slide (\texttt{main.pdf},
\texttt{main45.pdf}) và rà lỗi biên dịch; đối chiếu lần cuối mục lục sách với nội dung. \\
\bottomrule
\end{longtable}

\subsection{Sản phẩm bàn giao}
\begin{itemize}[leftmargin=1.2em]
    \item \texttt{DOCS/BOOK/} -- sách kỹ thuật hoàn chỉnh, 18 chương + lời giải.
    \item \texttt{DOCS/Slides67/main.pdf} -- slide tra cứu đầy đủ.
    \item \texttt{DOCS/Slides67/Slides45min/main45.pdf} -- slide báo cáo 45 phút.
    \item Bộ script + hình \texttt{matplotlib} sinh từ dữ liệu thật (\texttt{DOCS/BOOK/adc\_figures/real\_*.py}).
    \item Bảng số liệu benchmark SADGS chính thức (giai đoạn G2--G3).
\end{itemize}

\section{Hiện trạng tại thời điểm đề xuất}

Bảng dưới tóm tắt kết quả đợt rà soát G0, được thực hiện bằng cách đối chiếu song song từng
phần tài liệu với \texttt{SADGS/} qua nhiều lượt kiểm tra độc lập:

\begin{longtable}{@{}p{4.2cm} p{8.9cm}@{}}
\toprule
\textbf{Hạng mục} & \textbf{Kết quả} \\
\midrule
\endhead
Hình matplotlib bịa/không liên quan & Xoá khoảng 65 hình dùng số liệu \texttt{np.random}
hoặc công thức không tồn tại trong code; tạo mới khoảng 39 hình từ ảnh/checkpoint/log thật. \\
Cơ chế ``FastGS'' bịa & Gỡ bỏ/viết lại toàn bộ tuyên bố trình bày FastGS như baseline đã cài
đặt (395+ lần nhắc trong sách, 88+ lần trong slide) -- chỉ giữ lại như trích dẫn công trình
liên quan. \\
Cơ chế ``Importance score'' pixel-error & Đính chính toàn bộ slide/sách mô tả sai cơ chế
prune/clone thật (\texttt{accum\_view\_count > 0.5}, không phải điểm lỗi pixel). \\
Dead code bị mô tả như đang chạy & Chỉ rõ \texttt{final\_prune\_structgs},
\texttt{expand\_undersized\_gs} bị comment out trong \texttt{train.py}, không chạy ở cấu hình
mặc định. \\
Package \texttt{pipeline/} bịa & Viết lại chương 1--4 để tách rõ kiến trúc thật của
\texttt{SADGS/} (một entry point duy nhất, \texttt{train.py}) khỏi cấu trúc của fork ngoài
``SADGS-lite''. \\
Trích dẫn file:line sai & Sửa hàng loạt trích dẫn lệch dòng trong \texttt{16-loi-giai/} (do
\texttt{gaussian\_model.py} đã được mở rộng nhiều so với bản 3DGS gốc mà lời giải trích theo). \\
Lỗi biên dịch LaTeX & Sửa 2 lỗi khiến \texttt{main.tex} và \texttt{part7.tex} không build được
(comment \texttt{\%} lồng trong đối số brace; thiếu \texttt{[fragile]} cho frame chứa
\texttt{verbatim}) -- cả hai bộ slide hiện build thành công (339 trang và 58 trang). \\
\bottomrule
\end{longtable}

\section{Hạn chế và rủi ro đã biết}

\begin{itemize}[leftmargin=1.2em]
    \item \textbf{Số liệu benchmark hiện tại chưa phải của SADGS thật.} Bảng số liệu trong
    \texttt{DOCS/README.md} (Score 0.8579, PSNR 25.27\,dB, \ldots) là của pipeline
    ``FastGS-lite'' \emph{trước khi} tích hợp \texttt{densify\_and\_split\_structgs} et al. --
    đây là rủi ro lớn nhất nếu không thực hiện Giai đoạn G2, vì mọi so sánh định lượng
    ``SADGS nhanh hơn X\%'' hiện tại đều chưa có cơ sở.
    \item \textbf{Môi trường viết tài liệu không có \texttt{torch}/GPU.} Toàn bộ hình minh
    hoạ phải tái hiện công thức bằng \texttt{numpy} thuần -- rủi ro sai khác nhỏ so với bản
    CUDA/PyTorch thật nếu công thức bị tái hiện sai; mọi script cần trích dẫn \texttt{file:line}
    để người đọc tự đối chiếu lại.
    \item \textbf{Một số nhánh code không active có thể được kích hoạt lại trong tương lai}
    (ví dụ nếu \texttt{final\_prune\_structgs} được un-comment trong một bản cập nhật
    \texttt{SADGS/} sau này) -- tài liệu cần cơ chế theo dõi để không bị lỗi thời.
    \item \textbf{Chưa chạy lại toàn bộ 18--16-loi-giai/ và 18-sadgs-co-che-chuyen-sau/ với
    quy trình xác minh mới nhất} -- phần lớn đã được kiểm chứng trong đợt G0, nhưng nội dung
    phát sinh sau này cần đi qua checklist Mục~4.1 trước khi merge.
\end{itemize}

\section{Kết luận}

Dự án này không đề xuất một thuật toán mới, mà đề xuất một \textbf{kỷ luật tài liệu hoá}: mọi
tuyên bố về SADGS phải truy vết được tới mã nguồn hoặc dữ liệu thật. Đợt rà soát G0 (đã hoàn
thành phần lớn) cho thấy kỷ luật này là cần thiết -- một lượng lớn nội dung trước đó không
khớp với code đang chạy thật, dù không cố ý. Các giai đoạn còn lại (G1--G4), đặc biệt việc
chạy lại benchmark chính thức cho SADGS thật (G2), là điều kiện để tài liệu có thể được dùng
làm cơ sở so sánh định lượng đáng tin cậy.

\end{document}
